% !Mode:: "TeX:UTF-8"
%!TEX program  = xelatex
\documentclass[UTF8,zihao=-4]{ctexart}
\usepackage[a4paper,margin=2.4cm]{geometry}
\usepackage{amsmath,amssymb}
\usepackage{booktabs}
\usepackage{array}
\usepackage{enumitem}
\usepackage{xcolor}
\usepackage[hidelinks]{hyperref}

\setlength{\parindent}{2em}
\newcommand{\yuanduan}[1]{\par\noindent\textcolor{gray}{\small 【原文】}#1\par}
\newcommand{\xinduan}[1]{\par\noindent\textcolor{black}{\small 【修改后】}#1\par}

\title{\heiti 《药材烘干过程传热传质模型的建立与数值求解》论文审计与修改意见（v1）}
\author{}
\date{2026 年 9 月 12 日}

\begin{document}
\maketitle

\begin{center}
\textbf{说明}：本意见对应修改后的论文 \texttt{论文v1/论文.tex}（PDF 见 \texttt{论文v1/论文.pdf}）。
所有结论均与 \texttt{代码/problem1$\sim$4.py}、\texttt{表格/表1\_表2$\sim$表6.txt}、
\texttt{论文/code/sensitivity.json}、\texttt{论文/code/算法对比表.txt} 等计算记录逐项核对。
\end{center}

\section{审计结论总览}

\subsection{问题间模型的统一性与关联性（要求 1）}

经检查，四个问题的模型框架是统一的（问题一是问题二的常物性特例，问题三与问题二同模型，
问题四是问题三在材料坐标下的推广），但存在\textbf{两处表述与数据矛盾}：

\begin{enumerate}
\item 摘要声称问题二"与问题一的解衔接良好"，但表 1 与表 3 在 $t=1800\,\mathrm{s}$
（$=0.5\,\mathrm{h}$）处中心温度分别为 $33.5743\,{}^\circ$C 与 $32.1879\,{}^\circ$C，
相差约 $1.4\,{}^\circ$C；表 2 与表 4 的表面水分浓度分别为 $1.4613$ 与
$1.6483\,\mathrm{kg/kg}$，相差约 $12\%$。该表述与论文自身的数据矛盾（修改见 M1、M8）；
\item 问题四与问题三的对比分析断言"路径缩短占优"，但表 5 与表 6 显示问题四在约 42 h
之前中心水分浓度一直\textbf{高于}问题三（6 h 时 $1.7194$ 对 $1.0171$），后期才反超，
"路径缩短占优"的单一归因与数据趋势不符（修改见 M11）。
\end{enumerate}

\subsection{灵敏度分析是否成立（要求 2）}

灵敏度分析的数值结果与 \texttt{sensitivity.json} 完全一致，各项结论
（$h$ 不敏感、$D_0$ 与 $R$ 敏感、$C^{*}$ 最敏感、Arrhenius 温度效应约 $3\%/{}^\circ$C）
均成立，\textbf{不存在数据造假或计算错误}；但存在三处逻辑缺陷：

\begin{enumerate}
\item \textbf{摘要与正文结论矛盾}：摘要称"烘干时间对水分扩散系数与药材半径最敏感"，
而正文明确得出"目标含水率是最敏感因素"（$+47.04\%$），且 $C^{*}$ 属工艺要求而非物性参数，
二者应区分表述（修改见 M2）；
\item \textbf{敏感性数字不对称却写成对称区间}：$D_0\times0.8$ 使时间增加 $21.94\%$，
$D_0\times1.2$ 使时间缩短 $14.48\%$，原文"约 $\pm22\%$"与数据不符，应为"$+22\%/-14\%$"
（修改见 M14）；
\item \textbf{基准模型的可外推性未说明}：灵敏度分析以问题三（附录 3 物性）为基准，而论文的
最终烘干时间结论来自问题四（附录 4 物性），未说明问题三的灵敏度结论能否推广到问题四
（修改见 M14，补充基准选择理由与定性外推说明）。
\end{enumerate}

\subsection{内容缺陷检查（要求 3）}

\begin{enumerate}
\item \textbf{摘要}：已逐问覆盖问题一至四（原稿合格），仅两处表述需修正（M1、M2、M3）；
\item \textbf{假设}：缺两条必要假设——"恒温干燥阶段烘房温湿度恒定（取附件 1 平稳段稳定值）"、
"收缩仅沿径向、长度不变"（M5）；
\item \textbf{符号}：$T_{\mathrm{set}}$、$C_{\mathrm{set}}$、$C^{*}$、$D_0$、
$\lambda_T$、$\lambda_D$ 在正文与灵敏度表中使用但未列入符号表，$R$ 未注明问题四中随时间变化（M6）；
\item \textbf{数值方法介绍}：各数值方法均有介绍（合格），但问题一"表面虚拟节点"仅一句话带过、
未给出公式（M7）；且问题一稳定性参数 $\lambda_D\approx0.02$ 计算有误，按 $D_{\max}=7\times10^{-9}$
应为 $\lambda_D\le0.007$（M7）；问题四 $N=160$ 应表述为子区间数（M10）。
\end{enumerate}

\section{逐条修改意见}

以下每条给出\textbf{修改处}、\textbf{修改原因}、\textbf{原文}与\textbf{修改后文段}。

\subsection*{M1 摘要：问题二与问题一的衔接表述（与表 1$\sim$表 4 数据矛盾）}

\textbf{修改处}：摘要第二段末句。\\
\textbf{修改原因}：表 1$\sim$表 4 显示两问在同一时刻（$1800\,\mathrm{s}=0.5\,\mathrm{h}$）的
中心温度相差 $1.4\,{}^\circ$C、表面水分浓度相差 $0.19\,\mathrm{kg/kg}$，
"衔接良好"与数据矛盾，评审易被证伪；应改为说明两问的模型继承关系并指引读者到第 5 节的衔接说明。

\yuanduan{……给出 3 h 内的温度与水分浓度分布，并与问题一的解衔接良好。}
\xinduan{……给出 3 h 内的温度与水分浓度分布；该模型将问题一的常物性模型推广为变物性形式，
两问结果的差异与衔接在第 5 节说明。}

\subsection*{M2 摘要：灵敏度结论与正文矛盾（目标含水率）}

\textbf{修改处}：摘要末段灵敏度表述。\\
\textbf{修改原因}：正文灵敏度分析明确"目标含水率是最敏感因素"（$+47\%$），摘要只提
"扩散系数与半径最敏感"，属前后矛盾；且 $C^{*}$ 是工艺要求而非物性参数，应分层表述。

\yuanduan{灵敏度分析表明，烘干时间对水分扩散系数与药材半径最敏感，而对对流换热系数几乎不敏感，
说明干燥过程由内部扩散控制；据此给出降低烘房湿度、适当提高温度等工艺优化建议。}
\xinduan{灵敏度分析表明，物性参数中烘干时间对水分扩散系数与药材半径最敏感，工艺参数中
目标含水率的影响最大，而对流换热系数几乎不影响烘干时间，说明干燥过程由内部扩散控制；
据此给出降低烘房湿度、适当提高温度、合理设定目标含水率等工艺优化建议。}

\subsection*{M3 摘要：隐式格式效率表述（1/20 与算法对比表不符）}

\textbf{修改处}：摘要第三段。\\
\textbf{修改原因}：算法对比表（表 7）实测显式 $\Delta t=1\,\mathrm{s}$ 耗时约 20 s、
隐式 $\Delta t=60\,\mathrm{s}$ 耗时约 0.5 s，为 $1/40$ 而非 $1/20$，按实测数据修正。

\yuanduan{隐式格式无条件稳定，$\Delta t=60\,\mathrm{s}$ 时耗时仅需显式的约 1/20，……}
\xinduan{隐式格式无条件稳定，$\Delta t=60\,\mathrm{s}$ 时计算耗时约为显式格式的 1/40，……}

\subsection*{M4 问题分析：问题一"常物性"表述与模型统一性段落}

\textbf{修改处}：3 节"问题分析"之问题 1 段、以及"总体求解思路"之前。\\
\textbf{修改原因}：问题一的 $D=7\times10^{-9}\mathrm{e}^{-0.89/C}$ 随 $C$ 变化，
并非全部常物性，原文"物性按附录 2 取常数"不严谨；同时为直接回应"问题间模型应统一"的要求，
补写四问模型统一性说明段。

\yuanduan{\textbf{问题 1}（预热平衡阶段，30 min）：物性按附录 2 取常数，……可用显式有限差分法直接求解。}
\xinduan{\textbf{问题 1}（预热平衡阶段，30 min）：$\rho$、$c_p$、$k$、$h$、$h_m$ 按附录 2 取常数，
水分扩散系数按附录 2 经验公式随 $C$ 变化，……}

新增段落（插于"总体求解思路如图"之前）：\\[2pt]
\noindent\textcolor{black}{\small 【修改后】总体而言，四个问题在模型上层层递进而保持统一：问题一正是问题二
在常物性下的特例（守恒型扩散方程在物性取常数时退化为问题一的导热方程与 Fick 扩散方程）；
问题三与问题二采用完全相同的物理模型，仅将模拟时域由 3 h 延长至烘干结束；问题四是问题三在材料坐标下的
推广，收缩的全部影响归结为因子 $1/R^2(t)$。数值方法上，问题一为显式有限差分，问题二为守恒型显式有限体积
（物性变化时更稳健），问题三、四为隐式有限体积，各格式均基于同一径向网格与相同的边界处理，仅在时间推进
方式上按时间尺度需求选择显式或隐式。}

\subsection*{M5 模型假设：补充两条必要假设}

\textbf{修改处}：2.1 节假设 6、假设 7 之后。\\
\textbf{修改原因}：（1）问题四只利用附件 2 的半径收缩数据，长度变化被忽略，必须显式假设
"收缩仅沿径向、长度不变"；（2）问题二至四恒温阶段取 $50\,{}^\circ$C、$0.04986\,\mathrm{kg/kg}$，
这一取值依据应上升为模型假设。

\yuanduan{6. 失水收缩是均匀的：每个材料点保持其相对半径位置 $r/R(t)$ 不变，收缩规律由附件 2 给定；\\
7. 物性参数按附录 2$\sim$4 的常数或经验公式取值。}
\xinduan{6. 失水收缩是均匀的且仅沿径向发生（长度不变）：每个材料点保持其相对半径位置
$r/R(t)$ 不变，收缩规律由附件 2 给定；\\
7. 物性参数按附录 2$\sim$4 的常数或经验公式取值；\\
8. 恒温干燥阶段烘房温度与水分浓度保持恒定，取附件 1 平稳段的稳定值（$50\,{}^\circ$C、
$0.04986\,\mathrm{kg/kg}$）。}

\subsection*{M6 符号表：补充缺失符号}

\textbf{修改处}：2.2 节符号表。\\
\textbf{修改原因}：$T_{\mathrm{set}}$、$C_{\mathrm{set}}$、$C^{*}$、$D_0$、$\lambda_T$、$\lambda_D$
在正文、灵敏度表与图中使用但未定义；$R$ 未注明问题四中随时间收缩。

\yuanduan{\small $R$ 与 药材半径 与 m 一行；符号表末行为 $t_f$ 与 烘干结束时间 与 h。}
\xinduan{\small 将 "$R$" 行改为 "$R,\ R(t)$ \ 药材半径（问题四中随时间收缩，由附件 2 给定）"；
在 $T_{\mathrm{env}}$、$C_{\mathrm{env}}$ 后插入 $T_{\mathrm{set}}$（恒温干燥阶段烘房温度）、
$C_{\mathrm{set}}$（恒温干燥阶段烘房水分浓度）、$C^{*}$（目标水分浓度，烘干合格标准）；
在 $D$ 行后插入 $D_0$（扩散系数经验公式的前置因子）；
在 $\alpha$ 行后插入 $\lambda_T,\ \lambda_D$（温度、水分方程的稳定性参数）。}

\subsection*{M7 问题一：补充虚拟节点公式、修正 $\lambda_D$}

\textbf{修改处}：4.2 节显式离散段。\\
\textbf{修改原因}：（1）"表面节点引入虚拟节点处理"仅有文字而无公式，不符合
"数值方法应有相关介绍"的完整性要求，补出虚拟节点值公式并给离散式编号；
（2）$\lambda_D\approx0.02$ 计算有误：$D_{\max}=7\times10^{-9}\,\mathrm{m^2/s}$，
$\lambda_D=D\Delta t/\Delta r^2\le7\times10^{-3}$，原文偏大约 3 倍（不影响稳定性结论，但数字错误）。

\yuanduan{……表面节点 $r=R$ 引入虚拟节点并结合第三类边界条件式处理。……稳定性条件
$\lambda_T,\lambda_D\le1/2$ 显然满足（$\lambda_D\approx 0.02$），格式稳定。}
\xinduan{……表面节点 $r=R$ 引入虚拟节点 $T_{N+1}$、$C_{N+1}$，其值由第三类边界条件式的
中心差分近似给出：
\[
T_{N+1}^n = T_{N-1}^n - \frac{2\Delta r\,h}{k}\bigl(T_N^n-T_{\mathrm{env}}^n\bigr),\qquad
C_{N+1}^n = C_{N-1}^n - \frac{2\Delta r\,h_m}{D_N^n}\bigl(C_N^n-C_{\mathrm{env}}^n\bigr),
\]
再按与内部节点相同的差分模板计算表面值。……稳定性条件 $\lambda_T,\lambda_D\le1/2$ 显然满足
（$\lambda_T\approx0.169$，$\lambda_D=D\Delta t/\Delta r^2\le0.007$，因
$D\le7\times10^{-9}\,\mathrm{m^2/s}$），格式稳定。}

\subsection*{M8 问题二：新增"与问题一的衔接说明"、补充稳定性校核}

\textbf{修改处}：5.1 节末（阶段分界依据之后）与 5.2 节末。\\
\textbf{修改原因}：（1）用数据如实说明问题一、二在同一时刻的差异及其物理来源，并指出
"问题一是问题二的常物性特例、后文以问题二模型为基准"，修复原摘要"衔接良好"与表 1$\sim$表 4 的矛盾；
（2）原稿只给出初始 $\alpha$，未论证变系数下显式格式全程稳定，补充全程
$\lambda_T\le0.23$、$\lambda_D\le0.013$ 的校核；（3）$D$ 范围下限按实际计算值改为约 $5.6\times10^{-9}$。

\yuanduan{\small （5.1 节末）依据：附件 1 中烘房温度在约 1.92 h 首次达到 $49.9\,{}^\circ$C，……故以 2 h 为阶段分界。}
\xinduan{\small 【修改后】（承接其后）与问题一的衔接说明。问题二在同一时刻的结果与问题一存在小幅差异：
$t=1800\,\mathrm{s}$ 时，中心温度分别为 $33.5743\,{}^\circ$C（问题一）与 $32.1879\,{}^\circ$C（问题二），
表面水分浓度分别为 $1.4613$ 与 $1.6483\,\mathrm{kg/kg}$，而中心水分浓度分别为 $2.5500$ 与
$2.5499\,\mathrm{kg/kg}$、几乎一致。该差异源于附录 2（常物性，$k=0.36\,\mathrm{W/(m\cdot K)}$）与附录 3
（变物性，初始 $k=0.4830\,\mathrm{W/(m\cdot K)}$）对同一药材给出了不同的物性参数化，致使预热阶段的温升速率
与表层干燥速率略有不同；水分场在未受干燥前沿影响的内部区域则几乎完全重合。从模型结构看，问题一恰是问题二的
常物性特例（守恒型方程在物性取常数时退化为问题一的导热方程与 Fick 扩散方程），后续问题三、四
均以问题二的变物性模型为基准，保证了四个问题在模型框架上的统一。}

\noindent\textcolor{black}{\small 【修改后】5.2 节末："全程网格稳定性参数 $\lambda_T=\alpha\Delta t/\Delta r^2\le0.23$、
$\lambda_D=D\Delta t/\Delta r^2\le0.013$，均满足显式格式的稳定性条件 $\lambda\le1/2$。"
（原"全程水分扩散系数在 $5.56\times10^{-9}\sim$"改为"约 $5.6\times10^{-9}\sim$"。）}

\subsection*{M9 精度表述统一（0.03\% 与 0.02\%）}

\textbf{修改处}：7.3 节算法选择结论第二条。\\
\textbf{修改原因}：6.2 节给出"偏差约 0.02\%"，7.3 节却写"约 0.03\%"，同一结果两种说法，统一为 0.02\%。

\yuanduan{……效率提升近两个数量级且精度损失可忽略（约 0.03\%），配合 $O(N)$ 的 Thomas 追赶法……}
\xinduan{……计算耗时由显式的约 20 s 降至约 0.5 s（约 40 倍），精度损失可忽略（约 0.02\%），
配合 $O(N)$ 的 Thomas 追赶法……}

\subsection*{M10 问题四：网格表述、结束半径统一、边界系数说明}

\textbf{修改处}：7.2 节（问题四数值求解与结果）。\\
\textbf{修改原因}：（1）$N=160$ 在程序中是子区间数（161 个节点），"160 个径向节点"不准确；
（2）论文前文"收缩至约 1.198 cm"与后文"半径已收缩至 1.200 cm"不一致，统一表述；
（3）边界条件式右端的 $R(t)$ 易被误读为 $h$、$h_m$ 随半径变化，补充说明。

\yuanduan{取 $N=160$ 个径向节点、$\Delta t=5\,\mathrm{s}$……半径已收缩至 $1.200\,\mathrm{cm}$。}
\xinduan{取 $N=160$ 个径向子区间（$\Delta\xi=1/160$）、$\Delta t=5\,\mathrm{s}$……
半径已收缩至约 $1.20\,\mathrm{cm}$（附件 2 的末端值为 $1.198\,\mathrm{cm}$）。}

\noindent\textcolor{black}{\small 【修改后】（边界条件式后新增）问题四移动表面边界条件式右端出现的 $R(t)$ 由坐标变换产生，
$h$、$h_m$ 本身仍为常数（假设 3）。}

\subsection*{M11 问题四：与问题三对比分析的机理解释重写}

\textbf{修改处}：7.3 节（与问题三结果的对比分析）。\\
\textbf{修改原因}：原稿断言"路径缩短占优"，但表 5、表 6 显示问题四前期反而慢于问题三
（6 h 中心 $1.7194$ 对 $1.0171$），约 42 h 后才反超，单一归因与数据矛盾，属于明显逻辑谬误；
改为与数据吻合的三因素解释，并如实描述"先慢后快"的交叉现象。

\yuanduan{原因有二：一是半径由 2.0 cm 收缩到 1.2 cm，扩散路径因子 $R^2$ 减小到原来的 36\%，
尽管附录 4 的扩散系数比附录 3 小，但\textbf{路径缩短占优}；二是附录 4 的 $D$ 对水分的依赖较弱
（$\mathrm{e}^{-0.30/C}$ 对 $\mathrm{e}^{-0.45/C}$），干燥后期 $D$ 衰减慢，尾段"拖尾"现象减轻，
因此总时长更短。}
\xinduan{对比表 5 与表 6 可见一个值得注意的现象：在约 42 h 之前，
问题四的中心水分浓度反而\textbf{高于}问题三（如 6 h 时 $1.7194$ 对 $1.0171\,\mathrm{kg/kg}$），
约 42 h 之后才反超并提前完成。其机理有三：
\begin{itemize}
\item \textbf{前期物性不利}：附录 4 的扩散系数在高水分段约为附录 3 的 $1/5\sim1/4$，导热系数也较低
（初始 $k=0.264$ 对 $0.483\,\mathrm{W/(m\cdot K)}$），升温与干燥均更慢，故前期问题四慢于问题三；
\item \textbf{路径缩短的加速作用}：半径由 2.0 cm 收缩至约 1.2 cm，$R^2$ 减小到原来的约 36\%，
等效扩散系数 $D/R^2(t)$ 随之放大约 2.8 倍，这一加速作用随收缩过程逐步显现；
\item \textbf{末段拖尾大幅缩短}：附录 4 的 $D$ 对水分的依赖较弱（$\mathrm{e}^{-0.30/C}$ 对
$\mathrm{e}^{-0.45/C}$），干燥末段 $D$ 衰减慢，"拖尾"现象减轻，后期单位时间内的失水速率明显高于问题三。
\end{itemize}
后两个加速因素在干燥中后期超过前期物性不利造成的减速，最终使总时长缩短约 9.5\%。}

\subsection*{M12 算法比较：直接求逆的复杂度表述}

\textbf{修改处}：8.2 节（线性方程组求解器的比较）。\\
\textbf{修改原因}：普通 LU 分解对三对角矩阵是 $O(N)$（带状消元），本文对比的是通用（稠密）LU 分解
（LAPACK），注明"通用"以免评审质疑复杂度论断。

\yuanduan{比较 Thomas 追赶法与直接矩阵求逆（LU 分解）求解同一三对角系统的平均耗时……}
\xinduan{比较 Thomas 追赶法与直接矩阵求逆（通用 LU 分解，LAPACK 实现）求解同一三对角系统的平均耗时……}

\subsection*{M13 算法选择结论：效率提升幅度表述}

\textbf{修改处}：8.3 节第二条。\\
\textbf{修改原因}："效率提升近两个数量级"夸大（实测约 40 倍、步数约 60 倍），改为实测数据支持的表述。

\yuanduan{……隐式格式无条件稳定，$\Delta t=60\,\mathrm{s}$ 时仅需约 3400 步，
效率提升近两个数量级且精度损失可忽略……}
\xinduan{……隐式格式无条件稳定，$\Delta t=60\,\mathrm{s}$ 时仅需约 3400 步，
计算耗时由显式的约 20 s 降至约 0.5 s（约 40 倍），精度损失可忽略……}

\subsection*{M14 灵敏度分析：基准模型说明、不对称敏感性表述、可推广性说明}

\textbf{修改处}：9 节（模型灵敏度分析）四处。\\
\textbf{修改原因}：（1）补充为何以问题三为基准、结论能否推广到问题四（消除"灵敏度结论与最终结论
来自不同模型"的疑虑）；（2）$D_0$ 敏感性不对称（$+21.94\%/-14.48\%$），原文"约 $\pm22\%$"不准确；
（3）$h$ 扰动最大影响为 $0.06\%$，原文"不足 0.1\%"可更精确；（4）$C^{*}$ 属工艺要求，单独说明。

\yuanduan{\small （表 9 之前）基线 $t_f=56.52\,\mathrm{h}$，结果见表 9 与龙卷风图。}
\xinduan{\small 【修改后】（承接其后）选择问题三模型为基准的原因：问题二、三、四共享同一传热传质框架，
问题三的计算量最小（$\Delta t=10\,\mathrm{s}$ 时单次模拟约 3 s），便于大规模参数扫描；问题四与问题三
在模型结构上同构（仅物性与 $R(t)$ 不同），各参数的作用方向一致，因此问题三的灵敏度结论在定性上适用于
问题四，定量数值则以问题三为基准。}

\noindent\textcolor{black}{\small 【修改后】结论第一条："对流换热系数 $h$ 扰动 $\pm20\%$ 时烘干时间变化
不超过 $0.07\%$，而水分扩散系数前置因子 $D_0$ 的影响达到约 $+22\%/-14\%$，药材半径 $R$ 的影响约
$\pm20\%$。"}

\noindent\textcolor{black}{\small 【修改后】结论第五条："……目标从 0.15 收紧到 0.13，烘干时间激增 47\%，
因为干燥末段 $D\propto\mathrm{e}^{-0.45/C}$ 急剧衰减。需说明的是，$C^{*}$ 属于工艺要求而非物性参数；
制定烘干标准时应在品质允许范围内合理放宽目标含水率。"}

\noindent\textcolor{black}{\small 【修改后】结论末尾新增："上述定量结果均以问题三（附录 3 物性）为基准；
问题四由于收缩与附录 4 物性的差异，烘干时间数值不同，但各参数的相对重要性排序（$C^{*}$、$R$、$D_0$ 敏感，
$h$ 几乎不敏感）保持一致。"}

\subsection*{M15 结果表标题：冠以题目要求的"表 1$\sim$表 6"编号}

\textbf{修改处}：全文六个结果表标题及正文引用。\\
\textbf{修改原因}：题目要求"按表 1、表 2（…表 6）格式给出结果"，论文表格应直接对应题目编号，
便于评审对照。

\yuanduan{\small 表题原为"30 分钟内药材的温度（$^\circ$C）""3 小时内药材的水分浓度（kg/kg）"
"药材烘干过程的水分浓度（kg/kg，考虑收缩）"等。}
\xinduan{\small 依次改为"表 1　30 分钟内药材的温度""表 2　30 分钟内药材的水分浓度"
"表 3　3 小时内药材的温度""表 4　3 小时内药材的水分浓度""表 5　药材烘干过程的水分浓度"
"表 6　药材烘干过程的水分浓度（考虑收缩）"；并在正文引用处注明"（即题目要求的表 1、表 2）"等。}

\subsection*{M16 问题三：补充空间网格收敛性检验，烘干时间修正}

\textbf{修改处}：6.1 节末、6.2 节（收敛性与精度校核）、6.3 节结果与摘要。\\
\textbf{修改原因}：原两项校核（时间步长无关性、显隐式对比）均在同一套 $\Delta r=0.1\,\mathrm{cm}$（$N=20$）
网格上进行，只能证明时间离散收敛与实现一致，无法暴露空间离散误差。空间收敛序列实测为
$N=20/40/80/160/320\to56.52/57.05/57.32/57.43/57.50\,\mathrm{h}$，增量逐次减半，Richardson 外推
$57.57\,\mathrm{h}$，原值 $56.52\,\mathrm{h}$ 系统性偏低约 $1.8\%$。

\yuanduan{……烘干结束时间 $t_f\approx203460\,\mathrm{s}=56.52\,\mathrm{h}\approx2.35$ 天……}
\xinduan{6.2 节新增（1）空间网格收敛性（$N=20/40/80/160/320\to56.52/57.05/57.32/57.43/57.50\,\mathrm{h}$，
Richardson 外推 $57.57\,\mathrm{h}$，并解释“达标时刻是解的病态泛函”：末段中心水分以约 $0.0011\,\mathrm{kg/(kg\cdot h)}$
下降，$0.001\,\mathrm{kg/kg}$ 的空间误差折合约 $1\,\mathrm{h}$）；原时间步长校核与显隐式对比保留为（2）（3），
并注明（3）在同一 $\Delta r=0.1\,\mathrm{cm}$ 网格上进行；6.3 节改为“烘干结束时间取空间网格 Richardson 外推值
$t_f\approx57.57\,\mathrm{h}$（$N=320$ 直接解为 $207000\,\mathrm{s}=57.50\,\mathrm{h}$，相差约 4 min），约 2.40 天”。
摘要同步改为 $57.57\,\mathrm{h}$（约 2.40 天）并提及空间网格加密与外推。}

\subsection*{M17 问题三：主算例与结果文件改用细网格，灵敏度分析细网格重算}

\textbf{修改处}：problem3.py、sensitivity.py、表 5、result3.xlsx、表 9（灵敏度表）及图（龙卷风图、扫描图）。\\
\textbf{修改原因}：表 5 与灵敏度原值均来自 $N=20$ 粗网格，继承约 $1.8\%$ 的时间偏差；细网格 $N=320$ 下
表 5 数值变化 $\le0.001$（如 54 h 中心 $0.1528\to0.1538$），但烘干时间与各灵敏度绝对数值需同步更新。

\yuanduan{problem3.py 中 solve(10.0, snap\_at=10800.0)（$\Delta r=0.1$ cm）；sensitivity.py 中 dr=0.001。}
\xinduan{problem3.py 主算例改为 solve(10.0, dr=0.02/320)（$N=320$），result3.xlsx 与表 5 按插值取数重算
（烘干时间 $57.50\,\mathrm{h}$，末行中心 $0.1499$）；sensitivity.py 改为 dr=0.02/320 重跑，基线 $57.50\,\mathrm{h}$
（与外推值 $57.57\,\mathrm{h}$ 相差约 $0.1\%$），表 9 各扰动行、龙卷风图与扫描图同步更新。}

\subsection*{M18 问题四：主算例改用 N=320，烘干时间 51.10→51.12 h}

\textbf{修改处}：problem4.py、7.2 节、表 6、result4.xlsx。\\
\textbf{修改原因}：问题四原用 $N=160$（$51.10\,\mathrm{h}$），细网格 $N=320$ 为 $51.12\,\mathrm{h}$，影响小但应统一
校核标准。

\yuanduan{取 $N=160$ 个径向子区间……收敛性校核 $N=20/40/80/160\to50.83/51.00/51.07/51.10\,\mathrm{h}$；
$t_f\approx183960\,\mathrm{s}=51.10\,\mathrm{h}$。}
\xinduan{取 $N=320$ 个径向子区间（$\Delta\xi=1/320$）……收敛性校核 $N=20/40/80/160/320\to
50.83/51.00/51.07/51.10/51.12\,\mathrm{h}$；$t_f\approx184020\,\mathrm{s}=51.12\,\mathrm{h}$；表 6 末行时间改为 51.12。}

\subsection*{M19 新增物质守恒定量校核；算法对比表同网格说明；连带数字更新}

\textbf{修改处}：6.2 节新增（4）；7.1 节表 7 之后；摘要与 7.3 节对比分析。\\
\textbf{修改原因}：（1）全文“守恒”此前只有定性说法，补充定量校核：$N=320$ 下含水量减少
$\Delta W=7.6209\times10^{-4}\,\mathrm{m^3}$ 与表面累计对流通量 $F=7.6209\times10^{-4}\,\mathrm{m^3}$ 相对偏差小于 $0.1\%$；
（2）表 7 各算法在同一 $\Delta r=0.1\,\mathrm{cm}$ 网格上比较，横向结论不受影响但绝对时间含空间误差，需加注说明；
（3）摘要与对比分析随新结果更新：$56.52/51.10\to57.57/51.12\,\mathrm{h}$，缩短 $6.4\,\mathrm{h}$、约 $11.2\%$，
反超时刻由“约 42 h”更正为“约 40 h（36$\sim$42 h 之间）”。

\subsection*{M20 附录：新增问题一至问题四完整程序（附录 C$\sim$F）}

\textbf{修改处}：附录部分（原附录 A、B 之后新增四节）。\\
\textbf{修改原因}：按提交要求，将修改后（细网格、本机路径）的四个求解程序完整源码收入论文附录，
依次命名附录 C、D、E、F。

\yuanduan{附录原仅含“主要程序求解流程”（附录 A）与“程序与数据文件清单”（附录 B）。}
\xinduan{新增：附录 C 问题一求解程序（problem1.py）、附录 D 问题二求解程序（problem2.py）、
附录 E 问题三求解程序（problem3.py，$N=320$ 细网格）、附录 F 问题四求解程序（problem4.py，$N=320$）；
源码以 tcolorbox+listings 排版（带行号、自动断行，可跨页）；为控制篇幅，附录代码已删去全部注释与
文档字符串，算法语句与 \texttt{代码\textbackslash} 目录中文件完全一致。}

\subsection*{M21 表格与图片尺寸统一调整、表 1/表 2 合并排版}

\textbf{修改处}：全文表格与图片尺寸、表 1 与表 2 的浮动体结构、表 9 字号。\\
\textbf{修改原因}：（1）六张结果表（表 1$\sim$6）缩放比例由 0.8 倍正文宽度缩小为 0.62 倍正文宽度，
各图片宽度统一缩小（正文图由 0.92 倍正文宽度缩至 0.66 倍，图 7 由 1.0 缩至 0.7，图 10 由 0.72 缩至 0.55，
图 11 由 0.8 缩至 0.62，图 12 由 0.92 缩至 0.72），使图表与正文比例协调；（2）表 1 与表 2 合并进同一
浮动体，保证同页显示；（3）表 9 改用 \texttt{small} 字号并收紧列距，确保单页完整显示。
数据均未改动；编译校验：表 1$\sim$表 9 全部单页完整、无跨页，表 1/表 2 同页，图 1$\sim$图 12
页码严格递增。

\subsection*{M22 图片再次放大并统一比例}

\textbf{修改处}：全文图片宽度。\\
\textbf{修改原因}：此前图片缩放偏小，视觉上与表格（0.62 倍正文宽）不协调。现将两联图统一为
0.7 倍正文宽、图 7 为 0.78、图 10 为 0.62、图 11 为 0.68、图 12 为 0.78，使图片与表格比例接近。
编译校验：表 1$\sim$表 9 仍全部单页完整、无跨页，表 1/表 2 同页，图 1$\sim$图 12 页码仍严格递增。

\section{修改汇总表}

\begin{center}
\begin{tabular}{lll}
\toprule
编号 & 位置 & 修改类型 \\
\midrule
M1 & 摘要第 2 段 & 修正与表 1$\sim$4 矛盾的"衔接良好"表述 \\
M2 & 摘要第 3 段 & 修正与正文矛盾的灵敏度结论（补目标含水率、区分参数类别） \\
M3 & 摘要第 3 段 & 效率 $1/20\to1/40$（与算法对比表一致） \\
M4 & 问题分析 & 修正"常物性"表述；新增四问模型统一性段落 \\
M5 & 模型假设 & 假设 6 补充"仅径向收缩"；新增假设 8（恒温段恒定条件） \\
M6 & 符号表 & 补充 $T_{\mathrm{set}}$、$C_{\mathrm{set}}$、$C^{*}$、$D_0$、$\lambda_T$、$\lambda_D$；$R\to R,R(t)$ \\
M7 & 问题一 & 补虚拟节点公式与离散式编号；$\lambda_D\approx0.02\to\le0.007$ \\
M8 & 问题二 & 新增"与问题一的衔接说明"；补全程稳定性校核；$D$ 范围下限修正 \\
M9 & 7.3 节 & 精度损失表述统一为 0.02\% \\
M10 & 问题四 & $N=160$ 子区间；结束半径统一为约 1.20 cm（1.198）；边界系数说明 \\
M11 & 问题四对比分析 & 重写机理（"先慢后快"三因素解释，删除与数据矛盾的"路径缩短占优"） \\
M12 & 8.2 节 & "通用 LU 分解（LAPACK 实现）" \\
M13 & 8.3 节 & "近两个数量级"改为"约 40 倍（20 s$\to$0.5 s）" \\
M14 & 灵敏度分析 & 基准模型选择说明；$+22\%/-14\%$ 不对称修正；$C^{*}$ 类别说明；可推广性说明 \\
M15 & 全文表题 & 冠以题目要求的表 1$\sim$表 6 编号 \\
\midrule
M16 & 问题三 & 补空间收敛性检验；烘干时间 56.52$\to$57.57 h（Richardson 外推） \\
M17 & 问题三/灵敏度 & 主算例与 result3/表 5 改用 $N=320$；灵敏度细网格重跑（基线 57.50 h） \\
M18 & 问题四 & 主算例 $N=160\to320$；烘干时间 51.10$\to$51.12 h；表 6/result4 重算 \\
M19 & 校核/说明 & 新增物质守恒定量校核；表 7 同网格说明；摘要与对比分析连带更新（11.2\%） \\
M20 & 附录 & 新增问题一$\sim$四完整程序，依次为附录 C/D/E/F \\
M21 & 全文图表 & 尺寸统一缩小；表 1/表 2 合并同页；表 9 缩小字号保证单页完整 \\
M22 & 全文图片 & 图片再次放大（两联图 0.7，与表格比例接近），重验无跨页乱序 \\
\bottomrule
\end{tabular}
\end{center}

\section{第二轮修订（M16$\sim$M19）后的数据说明}

第二轮修订前，表 5、表 6、表 9 与烘干时间均为 $\Delta r=0.1\,\mathrm{cm}$（问题四为 $N=160$）网格值；
第二轮已按上述 M16$\sim$M19 全部改为 $N=320$ 细网格值并同步更新摘要、正文、表格、图与
result3/result4.xlsx。对应程序：problem3.py、problem4.py（主算例网格）与论文/code/sensitivity.py
（dr=0.02/320），重跑日志见 \texttt{sensitivity\_fine\_run.log}。表 1$\sim$表 4（问题一、二，3 h 内）与
表 7（算法对比）、表 8（线性求解器）未涉及烘干终点，保持原网格结果不变。

\end{document}
